% Generated draft: requires mathematical and source review.
\documentclass[11pt]{article}
\usepackage[margin=25mm]{geometry}
\usepackage{fontspec}
\setmainfont{Latin Modern Roman}
\usepackage{amsmath,amssymb,mathtools,unicode-math}
\setmathfont{Latin Modern Math}
\usepackage{enumitem,needspace,xurl,hyperref}
\setlength{\parindent}{0pt}
\setlength{\parskip}{5pt}
\setlength{\emergencystretch}{3em}
\newcommand{\N}{\mathbb{N}}
\newcommand{\Z}{\mathbb{Z}}
\newcommand{\Q}{\mathbb{Q}}
\newcommand{\R}{\mathbb{R}}
\newcommand{\C}{\mathbb{C}}
\newcommand{\F}{\mathbb{F}}
\newcommand{\A}{\mathbb{A}}
\newcommand{\PP}{\mathbb{P}}
\newcommand{\E}{\mathbb{E}}
\newcommand{\eps}{\varepsilon}
\newcommand{\dd}{\,\mathrm{d}}
\begin{document}
\section*{UnsolvedMath ID: 20002701}
\textbf{Source problem code:} catalogue entry
\subsection{Definitions and mathematical statement}
This is a misnumbered, page-truncated fragment of the final bullet of Problem 1.2, not a complete Problem 6.5. It concerns the dependence of nonnesting structures \(NN_W\) on crystallographic root systems and their enumerative correspondence with other Coxeter--Catalan objects. For a crystallographic root system with simple roots \(\alpha_i\), the usual positive-root order is \(\alpha\preceq\beta\) when \(\beta-\alpha\) is a nonnegative integer combination of the simple roots; nonnesting objects are antichains. The intended task is an appropriate noncrystallographic extension, but the truncated source supplies no complete axioms or single formal conjecture.

The aid reports later unique surrogate root posets for \(I_2(m)\) and \(H_3\) under six standard properties, and impossibility of satisfying all six for \(H_4\). The six properties, references, and exact uniqueness convention are not supplied. That report is not a theorem that no useful \(H_4\) nonnesting structure of any kind can exist.

Its explicit rank-two obstruction fixes a normalization and labeled rays. For \(m\ge3\), take unit positive normals
\[
 \rho_j=(\cos(j\pi/m),\sin(j\pi/m)),\qquad j=0,\ldots,m-1,
\]
with simple roots \(\rho_0,\rho_{m-1}\). Define the real-cone order by
\[
 \rho_i\preceq\rho_j\quad\Longleftrightarrow\quad
 \rho_j-\rho_i\in\R_{\ge0}\rho_0+\R_{\ge0}\rho_{m-1}.
\]
The aid asserts that its distinct incomparable pairs are exactly
\[
 \{\rho_j,\rho_{m-1-j}\},\qquad0\le j<m-1-j\le m-1.
\]
There are no larger antichains. Hence its antichain polynomial and refinement marking simple roots are
\[
 1+mt+\lfloor m/2\rfloor t^2,
\]
\[
 \sum_{A\text{ antichain}}s^{|A\cap\{\rho_0,\rho_{m-1}\}|}t^{|A|}
 =1+2st+(m-2)t+s^2t^2+(\lfloor m/2\rfloor-1)t^2.
\]
For \(m\ge4\), the resulting total \(1+m+\lfloor m/2\rfloor\) differs from the dihedral Catalan target \(m+2\).

The stronger labeled-poset obstruction is: for \(m\ge4\), any poset on these \(m\) labeled rays with both simple roots minimal and exactly \(m+2\) antichains cannot make the chamber reflection \(\rho_j\mapsto\rho_{m-1-j}\) an order automorphism. Indeed, the empty and singleton antichains leave room for only one nonsingleton antichain, necessarily the pair of simple roots. The other \(m-2\) roots must form a chain above them; every automorphism fixes that chain pointwise, whereas the chamber reflection moves interior labels when \(m\ge4\). The obstruction concerns this exact label action, not every abstract poset with the Catalan cardinality.

Unit normalization and replacing integer-root order by a real cone are material restrictions. In crystallographic even dihedral types, the usual unequal root lengths can break this labeled-ray symmetry; the obstruction does not refute the ordinary \(B_2\) or \(G_2\) root posets. Nor does it settle the general noncrystallographic extension programme in the truncated source.
\subsection{Short English statement}
The truncated source concerns noncrystallographic nonnesting structures. The aid computes a unit-normal dihedral real-cone order with the wrong Catalan count and shows that a Catalan-sized poset on those same labeled rays, with both simples minimal, must break chamber-swap symmetry. This is a scoped obstruction, not a general nonexistence theorem.
\subsection{Sources}
\begin{itemize}
\item Local registry, permanent UnsolvedMath ID \(20002701\), exact public code \texttt{catalogue entry}; canonical location \texttt{aim-probability-notes.json}, \texttt{notes[142]}.
\item AIM, \emph{Braid groups, clusters, and free probability}, 2005 workshop problem list, printed pages 4--5, final bullet of Problem 1.2 according to the supplied verification; reference \url{https://aimath.org/WWN/braidgroups/braidgroups.pdf}. The canonical extraction is labeled Problem 6.5 and ends mid-reference. No PDF was accessed here. No proposer, original remarks, or literature entry is supplied.
\item The aid reports six-property surrogate-root-poset results for \(I_2(m),H_3,H_4\) without authors, exact axioms, theorem numbers, or complete citations.
\item Supplied research attempt 1: \texttt{partial\_result}, run and registry status \texttt{partially\_solved}, rigor label \texttt{proved}, low novelty confidence; classification \texttt{PARTIAL-PROGRESS}, created and updated 14 August 2026. The aid is machine generated, not peer reviewed, and requires independent verification before citation.
\end{itemize}
\end{document}
